\section{Congruence and residue classes}\label{ch06:sec:residues}
Fix a positive integer $m$. For integers $a$ and $b$, we write
\[
a\equiv b\pmod m,
\]
and say that $a$ is \emph{congruent to $b$ modulo $m$}, when $m$ divides $a-b$, that is, when $a-b=mk$ for some integer $k$. (The word ``divides'' was defined in Section~\ref{ch05:sec:euclid}.) In words, $a$ and $b$ differ by a whole number of steps of size $m$. A clock gives an everyday example with $m=12$: the times 15:00 and 3:00 put the hour hand in the same place, because $15-3=12$. For $m=4$, we have $1\equiv5\pmod4$, since $1-5=-4=4\cdot(-1)$, and $1\equiv-3\pmod4$, since $1-(-3)=4=4\cdot1$. On the other hand, $1\not\equiv2\pmod4$, because $1-2=-1$ is not a multiple of four.

Congruence modulo $m$ is an equivalence relation on $\ZZ$. We check the three properties of Definition~\ref{ch06:def:equivalence}, for arbitrary integers $a$, $b$, $c$.
\begin{itemize}
\item \emph{Reflexivity.} $a-a=0=m\cdot0$, so $a\equiv a\pmod m$.
\item \emph{Symmetry.} If $a\equiv b\pmod m$, then $a-b=mk$ for an integer $k$. So $b-a=m(-k)$, and $-k$ is an integer, which gives $b\equiv a\pmod m$.
\item \emph{Transitivity.} If $a\equiv b$ and $b\equiv c\pmod m$, then $a-b=mk$ and $b-c=m\ell$ for integers $k$ and $\ell$. So
\[
a-c=(a-b)+(b-c)=mk+m\ell=m(k+\ell),
\]
and $k+\ell$ is an integer, which gives $a\equiv c\pmod m$.
\end{itemize}
Because of symmetry, it does not matter whether we test $a-b$ or $b-a$: one is a multiple of $m$ exactly when the other is.

\subsection*{Residue classes}
The equivalence class of an integer $a$ under this relation is called its \emph{residue class modulo $m$} and is written $[a]_m$. (\emph{Residue} is an old word for remainder.) It consists of all integers congruent to $a$. An integer $x$ satisfies $x\equiv a\pmod m$ exactly when $x-a=mk$ for some integer $k$, that is, when $x=a+mk$. So the class consists of all integers of the form $a+mk$ with $k$ an integer:
\[
[a]_m=\{x\in\ZZ:x\equiv a\pmod m\}=\{a+mk:k\in\ZZ\}.
\]
For modulus four, the class of $1$ is
\[
[1]_4=\{\ldots,-7,-3,1,5,9,\ldots\}:
\]
start at $1$ and move up or down in steps of four. Any member of a class can be used to name the class, and a member used in this way is called a \emph{representative} of the class. The integer $5$ is another representative of the same class. Starting at $5$ and moving in steps of four reaches exactly the same integers, so $[5]_4=[1]_4$, even though $5\ne1$. Keep three levels apart here: the integer $1$; the set $[1]_4$ of all integers congruent to it; and the collection of all residue classes modulo four, which we describe at the end of this section.

By contrast, the class $[2]_4=\{\ldots,-6,-2,2,6,10,\ldots\}$ has no member in common with $[1]_4$. Two residue classes are always either equal, like $[1]_4$ and $[5]_4$, or without a common member, like $[1]_4$ and $[2]_4$. They never overlap only partly, and the same is true for the classes of every equivalence relation, as we now prove.

\par\noindent\begin{minipage}{\linewidth}
\begin{proposition}\label{ch06:prop:classes}
Let $\sim$ be an equivalence relation on a set $A$, and let $x,y\in A$.
\begin{enumerate}
\item[(a)] $x\in[x]$.
\item[(b)] $[x]=[y]$ if and only if $x\sim y$.
\item[(c)] The classes $[x]$ and $[y]$ are either equal or disjoint.
\end{enumerate}
\end{proposition}
\end{minipage}\par

\begin{proof}
(a) By reflexivity, $x\sim x$, and this says that $x\in[x]$.

(b) First suppose $[x]=[y]$. By part (a), $x\in[x]$, so $x\in[y]$, which means $x\sim y$.

Conversely, suppose $x\sim y$. We prove two containments. Let $z\in[x]$, so that $z\sim x$. Together with $x\sim y$, transitivity gives $z\sim y$, so $z\in[y]$. This proves $[x]\subseteq[y]$. For the reverse containment, symmetry turns $x\sim y$ into $y\sim x$. Now let $z\in[y]$, so that $z\sim y$. Together with $y\sim x$, transitivity gives $z\sim x$, so $z\in[x]$. This proves $[y]\subseteq[x]$, and the two containments give $[x]=[y]$.

(c) Suppose $[x]$ and $[y]$ are not disjoint. Then some element $z$ belongs to both, so $z\sim x$ and $z\sim y$. Symmetry turns $z\sim x$ into $x\sim z$, and transitivity applied to $x\sim z$ and $z\sim y$ gives $x\sim y$. By part (b), $[x]=[y]$. So two classes that share even a single element are equal.
\end{proof}

Each of the three properties of an equivalence relation does a job in this proof. Reflexivity puts every element into its own class, so no element is left out. Symmetry and transitivity make sure that classes cannot overlap partly.

A \emph{partition} of a set $A$ is a collection of nonempty subsets of $A$, called the \emph{pieces}, such that every element of $A$ lies in exactly one piece. In other words, the pieces cover the whole set and no two different pieces overlap. Proposition~\ref{ch06:prop:classes} says that the equivalence classes of any equivalence relation form a partition. Indeed, each class $[x]$ is nonempty, because it contains $x$ by part (a). Every element $x$ lies in at least one class, namely $[x]$, again by part (a). And no element lies in two different classes, because by part (c) two classes that share an element are equal. The birth-month classes of Section~\ref{ch06:sec:differences} are an example: the pieces are the sets of students born in the same month.

For congruence, part (b) gives a test that we will use again and again:
\[
[a]_m=[b]_m\quad\text{exactly when}\quad a\equiv b\pmod m,\ \text{that is, when $m$ divides $a-b$.}
\]
For example, $[1]_4=[5]_4$ because $4$ divides $5-1$, while $[1]_4\ne[2]_4$ because $2-1=1$ is not a multiple of four.

\subsection*{A standard name for each class}
Every residue class modulo $m$ contains \emph{exactly one} of the integers $0,1,\ldots,m-1$.

\emph{At least one.} Let $a$ be any integer. Division with remainder (Section~\ref{ch05:sec:euclid}) gives integers $q$ and $r$ with $a=qm+r$ and $0\le r<m$. Then $r-a=m\cdot(-q)$, so $r\equiv a\pmod m$, and $r$ belongs to $[a]_m$. For example, $-7=(-2)\cdot4+1$, so $1$ is a member of $[-7]_4$ that lies in $\{0,1,2,3\}$.

\emph{At most one.} Suppose, for a contradiction, that two different integers from the list $0,1,\ldots,m-1$ both belong to $[a]_m$. Call the smaller one $r$ and the larger one $s$, so that $0\le r<s\le m-1$. Then $s\equiv a$ and $r\equiv a$. Symmetry turns the second into $a\equiv r$, and transitivity gives $s\equiv r\pmod m$, so $s-r=mk$ for some integer $k$. Since $r<s$, the number $s-r=mk$ is positive. As $m$ is positive, $k$ must be positive too: if $k\le0$, then $mk\le0$. A positive integer is at least $1$, so $k\ge1$ and $s-r=mk\ge m$. This contradicts $s-r\le m-1$, which holds because $s\le m-1$ and $r\ge0$. So a class cannot contain two different integers from the list.

It follows that there are exactly $m$ residue classes modulo $m$, namely
\[
[0]_m,\ [1]_m,\ \ldots,\ [m-1]_m.
\]
Every class is one of these: if $r$ is the member of $[a]_m$ in the list, then $r\equiv a$, so $[a]_m=[r]_m$ by the test above. And no two of them are equal: if $[r]_m=[s]_m$ for different $r$ and $s$ in the list, then this one class would contain both $r$ and $s$, by part (a) of Proposition~\ref{ch06:prop:classes}, and we have just seen that this is impossible.

The ``at least one'' argument also tells us which integer of the list lies in $[a]_m$: it is the remainder of $a$ on division by $m$, which Section~\ref{ch05:sec:euclid} writes as $a\bmod m$. This gives a second form of the test: \emph{two integers are congruent modulo $m$ exactly when they have the same remainder on division by $m$.} Indeed, if $a\equiv b\pmod m$, then $[a]_m=[b]_m$, and the one member of this class in the list is both $a\bmod m$ and $b\bmod m$. Conversely, if $a$ and $b$ have the same remainder $r$, then $a\equiv r$ and $b\equiv r$, so $a\equiv b$ by symmetry and transitivity. For example, $17$ and $-7$ are congruent modulo $4$, since $17-(-7)=24=4\cdot6$, and indeed $17\bmod4=1$ and $(-7)\bmod4=1$.

Notice that the word ``mod'' is used in two ways. The statement $17\equiv-7\pmod4$ relates two integers, and like any statement it is either true or false. The expression $17\bmod4$ is a single integer, the remainder $1$. So $17\equiv1\pmod4$ is true, while $17=1$ is false.

\pauseq{Which of the integers $0,1,\ldots,6$ belongs to the class $[-20]_7$? Name two other representatives of the same class.}{Division with remainder gives $-20=(-3)\cdot7+1$ with $0\le1<7$, so the answer is $1$, and $[-20]_7=[1]_7$. Any integer of the form $1+7k$ is another representative, for instance $8$ and $-6$. A common slip is to write $-20=(-2)\cdot7-6$ and report $-6$. That equation is true, but $-6$ is not in the list $0,1,\ldots,6$: division with remainder always produces a remainder between $0$ and $m-1$, even for a negative number.}

The set of all residue classes modulo $m$ is written $\ZZ/m\ZZ$, where $m\ZZ$ stands for the set of all multiples of $m$. Read the notation as ``the integers, with differences that are multiples of $m$ ignored.'' For example,
\[
\ZZ/4\ZZ=\{[0]_4,[1]_4,[2]_4,[3]_4\}
\]
is a set with four elements, and each of those elements is itself an infinite set of integers. This is the third of the three levels mentioned earlier. More generally, the set of all equivalence classes of an equivalence relation on a set $A$ is called the \emph{quotient} of $A$ by the relation, and $\ZZ/m\ZZ$ is the quotient of $\ZZ$ by congruence modulo $m$.
